% -*-portada.tex-*-
% Este fichero es parte de la plantilla LaTeX para
% la realización de Proyectos Final de Carrera, protejido
% bajo los términos de la licencia GFDL.
% Para más información, la licencia completa viene incluida en el
% fichero fdl-1.3.tex

% Fuente tomada del PFC 'libgann' de Javier Vázquez Púa

\begin{titlepage}

  \begin{center}

    \includegraphics[scale=0.35]{upv.png} \\
    
    \vspace{2.0cm}
    
    \LARGE{\textbf{ESCUELA SUPERIOR DE INGENIERÍA}} \\
    
    \vspace{1.0cm}
    
    \Large{\textbf{INGENIERÍA SUPERIOR EN TELECOMUNICACIÓN}} \\
    
    \vspace{3.0cm}
%mimateo: ¿seguro que quedamso en poner este título?  
    \Large{ARDUINO Y MIPS} \\
    
    \vspace{2.0cm}
    
    \Large{Lluís Satorre González} \\
  
    \vspace{0.5cm}

    \large{\today}
    
  \end{center}
\end{titlepage}
